\documentclass[aspectratio=169,11pt]{beamer}
\usepackage[utf8]{inputenc}
\usepackage[T1]{fontenc}
\usepackage[french]{babel}
\usepackage{helvet}
\renewcommand{\familydefault}{\sfdefault}
\usepackage{hyperref}
\definecolor{bleu}{HTML}{0000FF}
\definecolor{fond}{HTML}{F4F4EF}
\setbeamercolor{background canvas}{bg=fond}
\setbeamercolor{structure}{fg=bleu}
\setbeamercolor{title}{fg=bleu}
\setbeamercolor{frametitle}{fg=bleu,bg=fond}
\setbeamercolor{normal text}{fg=black}
\setbeamertemplate{navigation symbols}{}
\setbeamertemplate{itemize item}{\textcolor{bleu}{\tiny$\blacksquare$}}
\setbeamertemplate{footline}{\hfill\footnotesize\insertframenumber\hspace{1em}\vskip6pt}
\hypersetup{colorlinks=true,urlcolor=bleu,linkcolor=bleu}
\title{Union africaine}
\subtitle{Fonctionnement et données, vue d'ensemble}
\author{Veille institutionnelle}
\date{2026-11-03}

\begin{document}
\begin{frame}\titlepage\end{frame}

\section{Histoire et mission}
\begin{frame}{Histoire et mission}
\begin{itemize}
\item Succède à l'OUA (1963), créée à Durban en 2002, siège à Addis-Abeba
\item 55 États membres (à vérifier)
\item Objectifs : unité, intégration, paix et sécurité, développement
\item Acte constitutif : droit d'intervention en cas de crimes graves, rejet des changements anticonstitutionnels
\end{itemize}
\end{frame}

\begin{frame}{Histoire et mission : ressources à suivre}
\small
Veille hebdomadaire.
\begin{itemize}
\item \href{https://au.int/en/overview}{Vue d'ensemble de l'UA}\\{\footnotesize Page institutionnelle, publication : hebdomadaire}
\end{itemize}
\end{frame}

\section{Organes et gouvernance}
\begin{frame}{Organes et gouvernance (1/2)}
\begin{itemize}
\item Conférence des chefs d'État et de gouvernement : organe suprême, sommets annuels
\item Conseil exécutif des ministres, Comité des représentants permanents
\item Commission : secrétariat exécutif, présidée par un président élu
\item Conseil de paix et de sécurité, Parlement panafricain, Cour africaine des droits de l'homme et des peuples
\end{itemize}
\end{frame}

\begin{frame}{Organes et gouvernance (2/2)}
\begin{itemize}
\item ECOSOCC : organe consultatif de la société civile
\end{itemize}
\end{frame}

\begin{frame}{Organes et gouvernance : ressources à suivre}
\small
Veille hebdomadaire.
\begin{itemize}
\item \href{https://au.int/en/organs}{Organes de l'UA}\\{\footnotesize Page institutionnelle, publication : mensuelle}
\item \href{https://au.int/en/commission}{Commission de l'UA}\\{\footnotesize Page institutionnelle, publication : mensuelle}
\item \href{https://au.int/en/member_states/countryprofiles2}{Pays membres}\\{\footnotesize Liste, publication : annuelle}
\end{itemize}
\end{frame}

\section{Financement}
\begin{frame}{Financement}
\begin{itemize}
\item Contributions statutaires des États membres, barème par pays
\item Dépendance historique aux partenaires extérieurs pour les programmes
\item Décision de Kigali : prélèvement de 0,2 \% sur les importations éligibles pour financer l'UA (à vérifier)
\item Fonds pour la paix de l'UA
\end{itemize}
\end{frame}

\begin{frame}{Financement : ressources à suivre}
\small
Veille hebdomadaire.
\begin{itemize}
\item \href{https://au.int/en/documents}{Documents de l'UA}\\{\footnotesize Base documentaire, publication : hebdomadaire}
\end{itemize}
\end{frame}

\section{Paix et sécurité}
\begin{frame}{Paix et sécurité}
\begin{itemize}
\item Architecture africaine de paix et de sécurité (APSA)
\item Conseil de paix et de sécurité de 15 membres
\item Missions de soutien à la paix, force africaine en attente
\item Suspension des États après changement anticonstitutionnel
\end{itemize}
\end{frame}

\begin{frame}{Paix et sécurité : ressources à suivre}
\small
Veille hebdomadaire.
\begin{itemize}
\item \href{https://www.peaceau.org}{Conseil de paix et de sécurité}\\{\footnotesize Site spécialisé, publication : hebdomadaire}
\item \href{https://au.int/en/pressreleases}{Actualités de l'UA}\\{\footnotesize Médiathèque, publication : hebdomadaire}
\end{itemize}
\end{frame}

\section{Intégration économique}
\begin{frame}{Intégration économique}
\begin{itemize}
\item ZLECAf : zone de libre-échange continentale, en vigueur en 2019, échanges depuis 2021, secrétariat à Accra
\item Communautés économiques régionales reconnues par l'UA : huit (à vérifier)
\item AUDA-NEPAD : agence de développement de l'UA
\item Protocole sur la libre circulation des personnes
\end{itemize}
\end{frame}

\begin{frame}{Intégration économique : ressources à suivre}
\small
Veille hebdomadaire.
\begin{itemize}
\item \href{https://au-afcfta.org}{ZLECAf}\\{\footnotesize Site spécialisé, publication : hebdomadaire}
\item \href{https://au.int/en/recs}{Communautés économiques régionales}\\{\footnotesize Page institutionnelle, publication : trimestrielle}
\item \href{https://www.nepad.org}{AUDA-NEPAD}\\{\footnotesize Site spécialisé, publication : mensuelle}
\end{itemize}
\end{frame}

\section{Données}
\begin{frame}{Données}
\begin{itemize}
\item Statistiques harmonisées : stratégie SHaSA et Institut statistique de l'UA, STATAFRIC (à vérifier)
\item Suivi de l'Agenda 2063 par rapports continentaux
\item Données économiques africaines plutôt à la BAD et à la CEA
\item Documents officiels et décisions des sommets sur au.int
\end{itemize}
\end{frame}

\begin{frame}{Données : ressources à suivre}
\small
Veille hebdomadaire.
\begin{itemize}
\item \href{https://au.int/en/documents}{Documents de l'UA}\\{\footnotesize Base documentaire, publication : hebdomadaire}
\item \href{https://au.int/en/agenda2063/overview}{Agenda 2063}\\{\footnotesize Page stratégique, publication : mensuelle}
\item \href{https://au.int/en/treaties}{Traités de l'UA}\\{\footnotesize Base de traités, publication : trimestrielle}
\end{itemize}
\end{frame}

\subsection{au.int : documents et traités}
\begin{frame}{au.int : documents et traités}
\begin{itemize}
\item Décisions et déclarations des sommets
\item Traités et conventions avec état des ratifications
\item Rapports de la Commission
\end{itemize}
\end{frame}

\begin{frame}{au.int : documents et traités : ressources à suivre}
\small
Veille hebdomadaire.
\begin{itemize}
\item \href{https://au.int/en/documents}{Documents de l'UA}\\{\footnotesize Base documentaire, publication : hebdomadaire}
\item \href{https://au.int/en/treaties}{Traités de l'UA}\\{\footnotesize Base de traités, publication : trimestrielle}
\end{itemize}
\end{frame}

\subsection{Suivi de l'Agenda 2063}
\begin{frame}{Suivi de l'Agenda 2063}
\begin{itemize}
\item Rapports continentaux de mise en œuvre
\item Tableau de bord de l'Agenda 2063 (à vérifier)
\end{itemize}
\end{frame}

\begin{frame}{Suivi de l'Agenda 2063 : ressources à suivre}
\small
Veille hebdomadaire.
\begin{itemize}
\item \href{https://au.int/en/agenda2063/overview}{Agenda 2063}\\{\footnotesize Page stratégique, publication : mensuelle}
\end{itemize}
\end{frame}

\subsection{Médiathèque : actualités et communiqués}
\begin{frame}{Médiathèque : actualités et communiqués}
\begin{itemize}
\item Communiqués de presse de la Commission
\item Actualités des sommets et du Conseil de paix et de sécurité
\end{itemize}
\end{frame}

\begin{frame}{Médiathèque : actualités et communiqués : ressources à suivre}
\small
Veille hebdomadaire.
\begin{itemize}
\item \href{https://au.int/en/pressreleases}{Actualités de l'UA}\\{\footnotesize Médiathèque, publication : hebdomadaire}
\item \href{https://www.peaceau.org}{Conseil de paix et de sécurité}\\{\footnotesize Site spécialisé, publication : hebdomadaire}
\end{itemize}
\end{frame}

\section{Classifications}
\begin{frame}{Classifications}
\begin{itemize}
\item Cinq régions de l'UA : Nord, Ouest, Centre, Est, Sud
\item Huit communautés économiques régionales : UMA, CEN-SAD, COMESA, CEEAC, CEDEAO, IGAD, SADC, EAC
\item Chevauchements d'appartenance entre communautés
\item Statut des États : membres actifs ou suspendus
\end{itemize}
\end{frame}

\begin{frame}{Classifications : ressources à suivre}
\small
Veille hebdomadaire.
\begin{itemize}
\item \href{https://au.int/en/recs}{Communautés économiques régionales}\\{\footnotesize Page institutionnelle, publication : trimestrielle}
\item \href{https://au.int/en/member_states/countryprofiles2}{Pays membres}\\{\footnotesize Liste, publication : annuelle}
\end{itemize}
\end{frame}

\section{Évaluation et contrôle}
\begin{frame}{Évaluation et contrôle}
\begin{itemize}
\item MAEP : mécanisme africain d'évaluation par les pairs sur la gouvernance
\item Audit des comptes de la Commission
\item Suivi des décisions du sommet et rapports de la Commission
\end{itemize}
\end{frame}

\begin{frame}{Évaluation et contrôle : ressources à suivre}
\small
Veille hebdomadaire.
\begin{itemize}
\item \href{https://aprm-au.org}{MAEP}\\{\footnotesize Site spécialisé, publication : trimestrielle}
\end{itemize}
\end{frame}

\section{Plans stratégiques}
\begin{frame}{Plans stratégiques}
\begin{itemize}
\item Agenda 2063 : cadre stratégique sur cinquante ans
\item Plan décennal de mise en œuvre : 2014-2023, puis 2024-2033 (à vérifier)
\item Aspirations : prospérité, intégration, gouvernance, paix, identité, société, rôle mondial
\item Projets phares : ZLECAf, marché unique du transport aérien, passeport africain
\end{itemize}
\end{frame}

\begin{frame}{Plans stratégiques : ressources à suivre}
\small
Veille hebdomadaire.
\begin{itemize}
\item \href{https://au.int/en/agenda2063/overview}{Agenda 2063}\\{\footnotesize Page stratégique, publication : mensuelle}
\item \href{https://au.int/en/documents}{Documents de l'UA}\\{\footnotesize Base documentaire, publication : hebdomadaire}
\end{itemize}
\end{frame}

\section{Critiques et débats}
\begin{frame}{Critiques et débats}
\begin{itemize}
\item Écart entre décisions et application
\item Dépendance financière envers les partenaires extérieurs
\item Réponses aux changements anticonstitutionnels et aux conflits
\item Chevauchement des communautés régionales
\end{itemize}
\end{frame}

\begin{frame}{Critiques et débats : ressources à suivre}
\small
Veille hebdomadaire.
\begin{itemize}
\item \href{https://au.int/en/pressreleases}{Actualités de l'UA}\\{\footnotesize Médiathèque, publication : hebdomadaire}
\end{itemize}
\end{frame}

\end{document}
